% A.76: full root and DE proof/source/exact-certificate audits PASS.
\subsection{Cyclic-index-two closures with an odd abelian factor}

\begin{statement}{A.76}{Proved}{The whole cyclic-index-two class and its odd abelian extensions}
Let \(F\) be a finite two-group with a cyclic subgroup of
index at most two, and let \(B\) be a finite abelian group
of odd order. Every actual connected finite \'{e}tale cover
of \(E\setminus\{O\}\) whose geometric Galois closure is
a quotient of \(F\times B\) satisfies both assertions
of Task A.5 at every coefficient prime. Every actual
intermediate, every original induced cusp packet and every
permitted open derived subgroup is included.

The nonabelian possibilities for \(F\) are exactly the
dihedral, generalized-quaternion, semidihedral and modular
families of A.70, A.72, A.73 and A.75. The statement
includes the small orders explicitly below. It is a
preservation theorem for an odd abelian factor; it asserts
neither cofinality nor either universal assertion of A.45,
and does not cover a product of two nonabelian Sylow factors.
\end{statement}

\begin{proof}
\paragraph{An elementary classification, including the small orders.}
A cyclic \(F\) is abelian. Otherwise choose a cyclic
subgroup \(\langle r\rangle\) of index two and order
\(2^a\), and an element \(s\) outside it. This subgroup
is normal, and the actual extension has a presentation
\[
 r^{2^a}=1,\qquad srs^{-1}=r^k,\qquad s^2=r^h,
 \quad k^2\equiv1\pmod{2^a},\quad (k-1)h\equiv0\pmod{2^a}.
                                                        \tag{A.76.1}
\]
Indeed the listed relations reduce every word to one of
\(2^{a+1}\) words \(r^i s^j\), and the actual group
has that order. Changing \(s\) to \(r^b s\) replaces
\(h\) by \(h+(1+k)b\) modulo \(2^a\).

For \(a\ge3\), write \(n=2^{a-1}\). The only roots
of \(k^2=1\) modulo \(2^a\) are
\[
                         1,\ -1,\ 1+n,\ -1+n.
                                                        \tag{A.76.2}
\]
To see this, \(k\) is odd, the two even numbers \(k-1,k+1\)
have greatest common divisor two, and one of them has
exact two-valuation one. Their product being divisible
by \(2^a\), the other is divisible by \(2^{a-1}\).
This gives exactly the four residues in (A.76.2).
The remaining normalizations are exact:
\begin{itemize}
\item If \(k=1\), the group is abelian. If \(h\) is odd,
\(s\) generates a cyclic group of order \(2^{a+1}\).
If \(h\) is even, a change of \(s\) sets \(h=0\),
and the group is \(C_{2^a}\times C_2\).
\item If \(k=-1\), the consistency condition gives
\(h=0\) or \(h=n\). These are the dihedral and
quaternion presentations, respectively.
\item If \(k=n-1\), consistency gives \(h=0\) or
\(h=n\), because \(\gcd(n-2,2n)=2\). Changing
\(s\) by \(rs\) interchanges these values, since
\(1+k=n\). This gives the semidihedral presentation.
\item If \(k=n+1\), consistency says that \(h\)
is even. Since \(\gcd(n+2,2n)=2\), changing \(s\)
sets \(h=0\). This is the modular presentation.
\end{itemize}

For \(a=2\), the only roots are \(k=1,-1\); the
same first two cases give the abelian groups, \(D_8\)
and \(Q_8\). For \(a=1\), \(k=1\) and the groups
are \(C_4\) and \(C_2^2\). For \(a=0\), the group
has order at most two and is cyclic. Thus no small
nonabelian case is omitted, and the semidihedral and
modular families first enter when \(n\ge4\).
Every quotient again has a cyclic subgroup of index at
most two: take the image of the chosen cyclic subgroup.
These are group-theoretic statements about the actual
closure, not identifications of curves from their genera.

\paragraph{The actual connected coprime pullback.}
A subgroup of a direct product of finite groups of coprime
orders is the product of its intersections with those
factors. Indeed powers of \((x,y)\), with exponents
chosen by Bezout, give \((x,1)\) and \((1,y)\) in the
same subgroup. Apply this also to a normal quotient kernel.
The actual closure group therefore has the form
\(F_0\times B_0\), where \(F_0\) is a quotient of
\(F\) and \(B_0\) is odd abelian. The actual covering
subgroup is \(H_2\times H_B\).

The abelian cover corresponding to \(B_0/H_B\) is
unramified over the compact elliptic curve: the peripheral
commutator is killed by an abelian quotient. After choosing
an origin it is an actual elliptic isogeny
\(\eta:E'\to E\) of odd degree. Let \(Z\) be the
proper compactification of the actual \(H_2\)-intermediate
of the \(F_0\)-closure. Its field has two-power degree
over the function field of \(E\), while the isogeny
field is Galois of odd degree. They are linearly disjoint,
so their actual compositum is the field of the geometrically
connected fibre product
\[
                              Y=Z\times_E E'.          \tag{A.76.3}
\]
Proper smooth normalization is understood. The removed
points are exactly pairs of a cusp of \(Z\) and a
point of \(\ker\eta\). This identifies the actual
covering map as well as the field; it does not select
an unspecified component of a formal fibre product.

\paragraph{The geometric input retained from the four families.}
If \(Z\) has genus one, its pullback is an unramified
isogeny cover of the proper elliptic curve, so A.40 applies.
If \(F_0\) is abelian the same case holds. Otherwise
A.70, A.72, A.73 or A.75 applies to the actual cover
\(Z\), including every intermediate. Those proofs give
potential good Jacobian reduction of special two-rank zero,
actual two-primary torsion cusp differences, and hence a
stable marked model at two in which every cusp has a
distinct companion joined by rational components. The
marked endpoint components are included in this rational
path. This geometric property, not merely an already
known fixed-group equality, is the input needed here.

\paragraph{The map extends, and compact type is checked separately.}
Make one finite local extension defining the good models
and all marked points. Base the Abel map of \(Z\) at
a cusp sent to \(O\), and compose its relative extension
with the actual homomorphism \(J_Z\to E\). The
compact-type Abel-map theorem extends this morphism
across nodes and any inserted rational components.
\cite[Theorem 1.1 and Proposition 3.2]{fghDupuyRabinoff}
Homomorphisms between the good abelian varieties extend
over their abelian schemes. This constructs the extension
\(\mathcal Z\to\mathcal E\) of the actual generic map.

The odd isogeny extends to a finite \'{e}tale isogeny
of good elliptic models: its finite flat kernel has order
invertible in the residue field. Thus
\(\mathcal Y=\mathcal Z\times_{\mathcal E}\mathcal E'\)
is a proper flat semistable model of the actual connected
generic curve, finite \'{e}tale over \(\mathcal Z\).
Compact type does not follow from this pullback assertion
alone. Instead apply the nilpotent ramified-cover theorem
to the actual geometric Galois closure over \(E\),
whose group is \(F_0\times B_0\). Its only branch
section is \(O\), with \'{e}tale closure on the good
elliptic model. Lehr--Matignon Theorem 6.13 supplies
potential good reduction of the closure Jacobian.
\cite[Theorem 6.13]{fghLehrMatignon2005}
Its actual norm onto \(J_Y\) gives potential good
reduction there as well, so the semistable dual graph
of \(\mathcal Y\) is a tree. No special-rank-zero
assertion for \(J_Y\) is needed.

\paragraph{Every selected rational path lifts literally.}
Take a selected cusp above \(P\) on \(Z\) and a
distinct companion \(Q\). Let \(T\) be the connected
rational subtree joining their marked vertices in the
geometric special fibre. Each component of \(T\) maps
constantly to the good elliptic curve, and those constants
agree across the nodes. Since its endpoint sections map
to \(O\), the whole subtree maps to the special \(O\).
The fibre of the good \'{e}tale odd isogeny above \(O\)
is a discrete set of \(\deg\eta\) reduced points.
The pullback of \(T\) is therefore literally that many
disjoint copies of \(T\).

The selected marked section chooses one copy. The same
copy contains a section above \(Q\), with the same
kernel point of \(\eta\). Distinct prime-to-two kernel
sections have distinct reductions; the two generic marked
points are distinct because \(P\ne Q\). They are
joined entirely by rational components. Stabilization
can contract rational components but cannot introduce a
positive-genus component on this path. The separately
proved tree property makes this the unique joining path
required by A.42.

A.42 now proves all full fixed-inertia assertions, and
A.59 proves the normalizer assertions for the nilpotent
closure. The finite extensions define the algebraic data
and reduction only. Geometric paths are used up to inner
conjugacy; the established finite-field restriction argument
retains every original cusp packet and permitted open
derived subgroup. This is the actual preservation argument
of A.71/A.74 applied to the classified fourth family as well.
No fixed element is lifted through a quotient, and no
cofinality statement follows.

The elementary congruence and small-order checks are retained
in \path{code/a_cyclic_index_two_abelian.py}. The complete
classification and geometric preservation are proved above,
independently of the finite range of that certificate.
\end{proof}
